\section{\\Исследование процессов обнаружения, сопряжения и установления \\соединения в сетях Bluetooth}

\subsection{Цель работы}
Получение базовых теоретических знаний о механизмах работы сетей Bluetooth, а также практических навыков управления процессами поиска (inquiry), сопряжения (pairing) и установления соединения (connection) с использованием официального стека BlueZ в ОС семейства Linux.

\subsection{Задачи}
\begin{itemize}[nosep]
    \item Изучить теоретические основы процедур Inquiry, Page, Pairing и Connection в архитектуре Bluetooth BR/EDR.
    \item Ознакомиться с утилитами управления стеком BlueZ и анализатором трафика \texttt{btmon}.
    \item На практике выполнить сканирование окружающих Bluetooth-устройств, зафиксировать их адреса и классы устройств.
    \item Выполнить процедуру сопряжения с заданным устройством.
    \item Установить активное логическое соединение и проверить его состояние.
    \item Протоколировать этапы работы для оформления отчета.
\end{itemize}

\subsection{Теоретический материал}

\subsubsection{Процесс обнаружения устройств (Inquiry и Page)}
В сетях Bluetooth обнаружение устройств состоит из двух последовательных фаз: \textbf{Inquiry} (поиск) и \textbf{Page} (вызов).
\begin{itemize}[nosep]
    \item \textbf{Inquiry scan}: Устройства, находящиеся в режиме обнаружения, периодически прослушивают эфир на специальных частотах (Inquiry hopping sequence).
    \item \textbf{Inquiry}: Инициирующее устройство (Inquirer) рассылает ID-пакеты по частотной сетке поиска. Устройства в режиме сканирования отвечают пакетом \textbf{FHS (Frequency Hop Synchronization)}, который содержит 48-битный адрес устройства (BD\_ADDR), тип часов, класс устройства (CoD) и параметры синхронизации.
\end{itemize}

После получения FHS-пакетов инициатор формирует список найденных устройств. Для перехода к следующему этапу используется процедура \textbf{Page}, в ходе которой устройства синхронизируются по часам и устанавливают канал управления (LMP --- Link Manager Protocol).

\subsubsection{Процесс сопряжения (Pairing и Bonding)}
Сопряжение --- это процесс безопасного обмена ключами между двумя устройствами для установления защищённого канала.
\begin{itemize}[nosep]
    \item \textbf{Pairing} (сопряжение): процедура обмена аутентификационными ключами. В современных спецификациях (версия 2.1+) используется механизм \textbf{SSP (Secure Simple Pairing)}, основанный на алгоритме Диффи-Хеллмана (ECDH). SSP поддерживает несколько моделей ассоциации: Numeric Comparison, Just Works, Passkey Entry и Out of Band.
    \item \textbf{Bonding} (связывание): сохранение сгенерированного \textbf{Link Key} в энергонезависимой памяти устройств. После связывания последующие соединения могут устанавливаться автоматически без повторного прохождения процедуры аутентификации.
\end{itemize}

На уровне HCI (Host Controller Interface) процедура сопряжения реализуется через обмен специфическими пакетами команд и событий между хостом (ОС) и Bluetooth-контроллером. Ключевые этапы сопровождаются следующими пакетами:
\begin{itemize}[nosep]
    \item \textbf{IO Capability Request/Response} (HCI Event 0x0B): устройства обмениваются информацией о возможностях ввода-вывода для выбора оптимального метода ассоциации SSP.
    \item \textbf{User Confirmation Request} (HCI Event 0x0D) и \textbf{Reply} (HCI Command 0x0C): используются в режиме Numeric Comparison для подтверждения пользователем совпадения 6-значного кода на обоих устройствах.
    \item \textbf{Link Key Request} (HCI Event 0x17) / \textbf{Link Key Reply} (HCI Command 0x0B): контроллер запрашивает у хоста сохранённый ключ связи; при первом подключении хост возвращает отрицательный ответ, что инициирует генерацию нового ключа.
    \item \textbf{Link Key Notification} (HCI Event 0x17): контроллер передаёт хосту вновь сгенерированный Link Key для сохранения в базе данных сопряжённых устройств.
    \item \textbf{Encryption Change} (HCI Event 0x08): событие, подтверждающее активацию аппаратного шифрования на уровне Baseband после успешной аутентификации.
\end{itemize}

В случае успешного сопряжения активируется шифрование на уровне контроллера (Baseband), обеспечивающее конфиденциальность передаваемых данных. Анализ HCI-пакетов в утилитах \texttt{btmon} или Wireshark позволяет отладить процедуру сопряжения, определить выбранный метод SSP и проконтролировать корректность установки защищённого канала.

\subsubsection{Установление соединения (Connection)}
После успешного Page/Pairing на уровне хоста (Host) создаётся логический канал через подуровень \textbf{L2CAP (Logical Link Control and Adaptation Protocol)}. L2CAP обеспечивает мультиплексирование нескольких протоколов более высокого уровня над одним физическим каналом.
Для установления соединения используются профили (например, \textbf{RFCOMM} для эмуляции последовательного порта, \textbf{A2DP} для аудио, \textbf{HID} для ввода). Клиент запрашивает соединение у серверного процесса, указывая PSM (Protocol/Service Multiplexer). После установки L2CAP-соединения устройства могут обмениваться полезными данными в рамках выбранного профиля.

\subsubsection{Стек BlueZ и основные инструменты управления}
BlueZ --- официальный стек протоколов Bluetooth для ядра Linux. Основной пользовательский интерфейс взаимодействия со стеком предоставляется через интерактивную утилиту \texttt{bluetoothctl}. Управление осуществляется через D-Bus API демона \texttt{bluetoothd}.

\begin{table}[htbp]
\centering
\caption{Основные команды утилиты \texttt{bluetoothctl}}
\label{tab:bluetoothctl}
\begin{tabularx}{\textwidth}{lX}
\toprule
\textbf{Команда} & \textbf{Описание} \\
\midrule
\texttt{power on} & Включение адаптера Bluetooth \\
\texttt{scan on} & Запуск режима сканирования (Inquiry) \\
\texttt{scan off} & Остановка сканирования \\
\texttt{devices} & Вывод списка обнаруженных или сопряжённых устройств \\
\texttt{pair <BD\_ADDR>} & Запуск процедуры сопряжения с указанным устройством \\
\texttt{trust <BD\_ADDR>} & Добавление устройства в список доверенных (автосопряжение) \\
\texttt{connect <BD\_ADDR>} & Установление логического соединения \\
\texttt{info <BD\_ADDR>} & Вывод подробной информации об устройстве (ключи, профили, статус) \\
\texttt{quit} / \texttt{exit} & Завершение сеанса работы с утилитой \\
\bottomrule
\end{tabularx}
\end{table}

Для мониторинга низкоуровневого трафика (HCI-команды, события, ACL-пакеты) в пакете BlueZ предусмотрен анализатор \texttt{btmon}, аналогичный по назначению утилите \texttt{tcpdump} в IP-сетях.

\subsection{Порядок выполнения лабораторной работы}
\textbf{Примечание.} Для выполнения работы требуется наличие активного Bluetooth-адаптера на лабораторном ПК. Все команды, кроме \texttt{info} и \texttt{devices}, требуют прав суперпользователя или работы в группе \texttt{bluetooth}.

\begin{enumerate}[nosep, leftmargin=*]
    \item Запустите терминал и откройте утилиту управления стеком:
    \begin{lstlisting}
user@host:~$ sudo bluetoothctl
    \end{lstlisting}
    В результате откроется интерактивная оболочка с приглашением \texttt{[bluetooth]\#}.

    \item Проверьте состояние адаптера и включите его при необходимости:
    \begin{lstlisting}
[bluetooth]# show
[bluetooth]# power on
    \end{lstlisting}
    Зафиксируйте вывод команды \texttt{show} (адрес адаптера, имя, текущее состояние).

    \item Запустите процесс сканирования устройств (Inquiry):
    \begin{lstlisting}
[bluetooth]# scan on
    \end{lstlisting}
    Подождите 10--15 секунд. В терминале будут появляться сообщения вида \lstinline|Device XX:XX:XX:XX:XX:XX|. После получения достаточного количества устройств выполните:
    \begin{lstlisting}
[bluetooth]# scan off
    \end{lstlisting}

    \item Выберите одно из обнаруженных устройств для дальнейшей работы. Выведите список найденных устройств и запомните его BD\_ADDR:
    \begin{lstlisting}
[bluetooth]# devices
    \end{lstlisting}

    \item Запустите параллельный терминал и начните захват низкоуровневого трафика с помощью \texttt{btmon}:
    \begin{lstlisting}
user@host:~$ sudo btmon -w bt_trace.log
    \end{lstlisting}
    Утилита будет записывать сырые HCI-события в файл \texttt{bt\_trace.log} в бинарном формате.

    \item В первом терминале инициируйте процедуру сопряжения (Pairing):
    \begin{lstlisting}
[bluetooth]# pair <BD_ADDR>
    \end{lstlisting}
    Следуйте инструкциям утилиты (при необходимости подтвердите PIN-код или код числового сравнения на обоих устройствах). Дождитесь сообщения \texttt{Pairing successful}.

    \item Установите логическое соединение (Connection):
    \begin{lstlisting}
[bluetooth]# connect <BD_ADDR>
    \end{lstlisting}
    Убедитесь в выводе \texttt{Connection successful}. Выполните команду \texttt{info <BD\_ADDR>} и зафиксируйте статус соединения, список доступных UUID и состояние шифрования.

    \item Завершите работу с \texttt{bluetoothctl} и остановите \texttt{btmon} (нажатие \texttt{Ctrl-C}). Просмотрите захваченный журнал:
    \begin{lstlisting}
user@host:~$ btmon -r bt_trace.log
    \end{lstlisting}
    Найдите в логе этапы: \texttt{Inquiry}, \texttt{Create Connection} (Page), \texttt{User Confirmation Request} / \texttt{Link Key}, \texttt{Encryption Change}.

    \item Выведите утилиту из интерактивного режима:
    \begin{lstlisting}
[bluetooth]# disconnect <BD_ADDR>
[bluetooth]# quit
    \end{lstlisting}
\end{enumerate}

\subsection{Содержание отчета}
\begin{enumerate}[nosep, leftmargin=*]
    \item Заголовок согласно приложению.
    \item Цель работы.
    \item Задание согласно варианту.
    \item Листинги результатов работы:
    \begin{enumerate}[label=\alph*), nosep]
        \item вывод команды \texttt{show} до и после включения адаптера;
        \item список обнаруженных устройств (\texttt{devices}) с указанием BD\_ADDR и класса устройства;
        \item вывод команд \texttt{pair}, \texttt{connect}, \texttt{info} с комментариями к каждому этапу;
        \item фрагменты лога \texttt{btmon}, иллюстрирующие пакеты Inquiry, Page и установку шифрования;
        \item скриншоты или текстовые подтверждения статуса соединения.
    \end{enumerate}
    \item Ответы на контрольные вопросы.
    \item Краткие выводы по работе.
\end{enumerate}

\textbf{Примечание.} Перед каждым листингом должна быть указана выполненная команда. Вывод \texttt{btmon} рекомендуется прилагать отдельным файлом или в виде вырезок, подтверждающих прохождение ключевых фаз протокола.

\section*{Контрольные вопросы}
\begin{enumerate}[nosep, leftmargin=*]
    \item Опишите различие между фазами Inquiry и Page. Какие пакеты используются на каждом этапе?
    \item Что такое FHS-пакет и какую информацию он передаёт инициатору поиска?
    \item В чём заключается разница между процессами Pairing и Bonding в спецификации Bluetooth?
    \item Какие модели ассоциации предусмотрены в механизме Secure Simple Pairing (SSP) и в каких сценариях они применяются?
    \item Какую роль играет подуровень L2CAP при установлении соединения между хостами?
    \item Как с помощью утилиты \texttt{btmon} отфильтровать только события, связанные с шифрованием канала?
    \item Почему для работы с адаптером в BlueZ часто требуются права суперпользователя или добавление пользователя в группу \texttt{bluetooth}?
\end{enumerate}
